\section{Counting and enumeration tools}
\label{sec:preliminaries}

Both trials use the same three facts about an arbitrary fixed input.
The number of distinct sums changes little when one index is replaced;
prescribed-cardinality choices have accurate entropy estimates; and a
random prime rarely identifies two fixed unequal integer weights.
We prove these facts first. We then give the small-state streams that
turn stored lists into much larger generated lists.

All logarithms are base two unless written $\ln$. Every finite decimal
constant is an exact rational. Polynomial factors below have absolute
degrees independent of the input-bit promise. We use the parameters
\begin{equation}
\label{eq:parameters}
\begin{gathered}
 u=n+b+2,\qquad m=8\lfloor3n/100\rfloor,\qquad
 s=m/2,\qquad g=m/4,\\
 C=2^{\lfloor .199n\rfloor},\qquad
 B=2^{\lfloor .19n\rfloor},\qquad h_*=h(1/4),
\end{gathered}
\end{equation}
where $h(x)=-x\log x-(1-x)\log(1-x)$ has its continuous endpoint
values. Blocks used for compression have size $s$; the representation
trial will use blocks of size $m$ and their halves and quarters.
The capacity $C$ counts distinct dictionary entries, whereas $B$
counts buffered occurrences in a stream. An entry can have many
occurrences, so these capacities have different roles.

The constructions below concern the \emph{main branch}, defined by
$n\ge n_0$ and
\begin{equation}
\label{eq:guard}
 u\le 2^{\lfloor n/10^9\rfloor},
\end{equation}
for one sufficiently large absolute integer $n_0$ fixed in the
program. Outside this branch the program uses exact subset
enumeration. Section~\ref{sec:implementation} gives the finite
implementation of this test, the repetition schedule, and the operation
cap. Each eventual fixed polynomial-bit class enters the main branch.

\subsection{How many distinct sums does a block have?}

For a domain $V\subseteq\{1,\ldots,n\}$, let $N(V)$ count the
\emph{distinct integer values} of its subset sums, and define
\[
 \defect(V)=|V|-\log N(V).
\]
We call this its deficiency. In particular, $N(\varnothing)=1$ and
$\defect(\varnothing)=0$. Here an actual weight is an integer sum,
as opposed to a residue; masks always encode sets of indices, even
when some input values agree. The next lemma applies also to signed
weights, which arise later after a sign transformation.

\begin{lemma}[Deficiency]
\label{lem:deficiency}
For arbitrary fixed weights on the $n$ indices, deficiency is
nonnegative, invariant under sign changes of individual weights,
inclusion-monotone, and superadditive on disjoint domains. Replacing
one index in a domain of fixed size changes its deficiency by at most
one. If $V$ is uniform among the size-$v$ domains, $0\le v\le n$,
then
\begin{equation}
\label{eq:deficiency-variance}
 \operatorname{Var}(\defect(V))\le v.
\end{equation}
\end{lemma}

\begin{proof}
There are $2^{|V|}$ subsets, so $N(V)\le2^{|V|}$.
Adding one index multiplies the number of sums by a factor between
one and two. Its addition therefore increases deficiency by a quantity
in $[0,1]$, proving inclusion monotonicity. For disjoint $V,W$,
every sum on their union is a sum of one value from each domain, so
$N(V\cup W)\le N(V)N(W)$. Taking logarithms gives
$\defect(V\cup W)\ge\defect(V)+\defect(W)$.

For sign invariance, let $F\subseteq V$ be the negated indices.
Toggling $F$ in a chosen subset is a bijection on all subsets of $V$.
The new weight of a subset is the old weight of its toggle minus the
old weight of $F$, so the set of sum values is merely translated.
Finally, two equal-sized domains differing in one index are extensions
of their common part. Each extension increases deficiency by a
quantity in $[0,1]$, which proves the replacement bound.

For the variance, expose a uniform ordered size-$v$ sample without
replacement. Let $M_r$ be the conditional expectation of the final
deficiency after the first $r$ positions, $0\le r\le v$.
Given the first $r-1$ positions, couple the remaining completions for
two possible next indices by transposing those indices. The final
domains then agree or differ in one index. Their deficiencies, and
hence their conditional means, differ by at most one. Thus all
possible values of $M_r$ given the past lie in an interval of length
at most one, and $|M_r-M_{r-1}|\le1$. Orthogonality of martingale
differences now gives
\[
 \operatorname{Var}(\defect(V))
 =\sum_{r=1}^v\mathbb E(M_r-M_{r-1})^2\le v.
\]
The empty-sample case is immediate.
\end{proof}

Consequently, for any fixed $\epsilon>0$, a deviation of
$\epsilon n$ from the size-specific mean has probability $O(1/n)$,
by Chebyshev's inequality. This concentration concerns a random
\emph{domain} in a fixed input; it assumes no distribution of the
numerical input values.

\subsection{Exact entropy estimates}

Cardinality restrictions on a block reduce its subset count from
$2^v$ to a binomial coefficient. We will need the exact exponential
rate, including in a later estimate with no exponential slack.

\begin{lemma}[Elementary entropy bounds]
\label{lem:entropy}
For integers $0\le k\le v$,
\[
 \frac{2^{v h(k/v)}}{v+1}\le \binom vk\le 2^{v h(k/v)}.
\]
For $v=k=0$, interpret $v h(k/v)$ as zero. Entropy is concave,
symmetric about $1/2$, and increasing on $[0,1/2]$. Moreover,
\begin{equation}
\label{eq:entropy-numbers}
 h_*<.8114,\qquad h(1/3)<14/15,\qquad h(1/6)>5/8.
\end{equation}
\end{lemma}

\begin{proof}
For $0<k<v$, consider the binomial distribution with success
probability $k/v$. The probability of exactly $k$ successes is
$\binom vk 2^{-v h(k/v)}$. It is at most one and at least
$1/(v+1)$: adjacent probability ratios show that $k$ is a mode,
and there are $v+1$ possible outcomes. This gives both estimates.
The cases $k=0,v$ and $v=0$ are immediate. Differentiation gives
the asserted monotonicity and concavity; symmetry follows from the
formula for $h$.

For the strict numerical inequalities, first
$h_*=2-\tfrac34\log3$. We have
$3^{12}/2^{19}>1+1/74$, whose fourth power exceeds
$1+4/74=39/37>256/243$. Hence $3^{53}>2^{84}$ and
$\log3>84/53>1.5848$, giving $h_*<.8114$.
Next, $h(1/3)=\log3-2/3$ and $3^5<2^8$, so
$h(1/3)<8/5-2/3=14/15$.
Finally, $6^2>2^5$ and $(6/5)^4>2$ imply
\[
 h(1/6)=\tfrac16\log6+\tfrac56\log(6/5)
 >\tfrac16\cdot\tfrac52+\tfrac56\cdot\tfrac14=\tfrac58.
\]
Thus none of the three bounds relies on floating-point optimization.
\end{proof}

\subsection{Random primes and the sampling interface}

The modular filters must control collisions of actual weights, even
when their bit length grows polynomially with $n$. We use the
following elementary estimate with its full bit-length dependence.

\begin{lemma}[Prime factors]
\label{lem:prime}
Let $4\le\ell\le n$ be an integer. The set
\[
 \mathcal P_\ell=
 \{p\text{ prime}:p<2^\ell,\ 4\ell p\ge2^\ell\}
\]
has at least $2^\ell/(4\ell)$ elements. If $p$ is uniform on this
set and $\Delta$ is a fixed nonzero integer with
$|\Delta|<2^{2u}$, then
\[
 \Pr[p\mid\Delta]\le 8u^2 2^{-\ell}.
\]
In particular, this estimate applies to every nonzero actual-weight
difference used by the two trials.
\end{lemma}

\begin{proof}
Put $X=2^\ell$. The central binomial coefficient obeys
$\binom{X}{X/2}\ge2^X/(X+1)$. In its factorial formula, the
exponent of a prime $p$ is
\[
 \sum_{r\ge1}
 \left(\left\lfloor\frac X{p^r}\right\rfloor
      -2\left\lfloor\frac {X/2}{p^r}\right\rfloor\right).
\]
Each summand is zero or one, and at most
$\lfloor\log_p X\rfloor$ of them are nonzero. Thus each prime
contributes a factor at most $X$ to the binomial coefficient. If
$\pi(X)$ is the number of primes at most $X$, taking logarithms gives
\[
 \ell\pi(X)\ge X-\log(X+1)\ge X/2
 \qquad(X\ge16).
\]
Remove primes below $X/(4\ell)$. There are at most $X/(4\ell)$
integers in that range, so at least $X/(4\ell)$ primes remain.
The upper endpoint $X$ is composite, proving the cardinality bound.

An integer $\Delta$ with $0<|\Delta|<2^{2u}$ has at most $2u$
distinct prime divisors: the product of $r$ distinct primes is at
least $2^r$. Dividing by the number of available primes and using
$\ell\le n\le u$ gives the displayed collision bound.
Actual weights and their relevant differences have magnitude
$O(n2^b)$; the constructions sum each indexed weight at most twice,
even before rejecting an overlapping tuple. They therefore meet
$|\Delta|<2^{2u}$ for sufficiently large main-branch inputs.
\end{proof}

We analyze uniform independent draws, conditional on each specified
finite range. The implementation uses at most $n^3$ attempts for each
draw and fails the current trial if they all fail. The precise
interface, proved in Lemma~\ref{lem:samplers}, is as follows. A draw
from $\{0,\ldots,D-1\}$ has the uniform law on success and couples
to an ideal draw with failure probability at most $2^{-n^3}$;
$D=1$ needs no randomness. A draw from $\mathcal P_\ell$ has its
uniform law on success and coupling failure at most $e^{-n^2/4}$.
These conditional bounds remain valid for ranges determined by past
choices. All requested prime parameters satisfy $\ell\le.258n$;
bounded trial division therefore costs at most
$\poly(n)2^{.129n}$ per prime search, including all its attempts.
The final section proves the finite-word implementation and the
aggregate coupling estimate. The trials below use no unbounded
sampling operation in their resource guarantees.
